\section{Introduction}
Contact-rich manipulation tasks, such as wiping and assembly, are ubiquitous in daily human activities. Their successful execution requires not only precise geometric alignment, but also accurate perception and fine-grained control of contact forces, frictional changes, and subtle transitions in contact state~\cite{zhao2024tac, cui2025vi, zhao2025tacman}. Since such information is difficult to infer reliably from vision alone, recent studies~\cite{bi2025vla, xue2025reactive, chen2025multi} have incorporated tactile sensing into imitation learning frameworks to improve visuomotor policies for contact-rich manipulation.

However, despite the growing adoption of tactile sensing in contact-rich manipulation, the broader visuo-tactile manipulation community still faces substantial limitations in both data and methodology. 
(1) On the data side, existing public visuo-tactile manipulation datasets~\cite{bi2025vla, yu2025demonstrating, bu2025agibot} remain limited in both the scale of fully aligned vision-tactile-action demonstrations and the diversity of contact-rich tasks. In particular, the amount of synchronized visual and high-frequency tactile data is still insufficient for learning generalizable visuo-tactile representations and modeling contact dynamics. In addition, the covered scenarios, tasks, and interaction patterns are narrow, making it difficult to capture the broad spectrum of contact mechanisms required in real-world manipulation.
(2) On the methodology side, existing visuo-tactile manipulation approaches remain limited in both representation and control. In most prior works~\cite{huang20243d, wu2025canonical}, tactile signals are incorporated merely as auxiliary observations to the policy network, primarily for contact-state recognition or compensating for visual occlusion, rather than being used to model the evolution of contact dynamics explicitly. At the control level, these methods typically rely on action chunking to generate short action sequences that are then executed in an open-loop manner. Such a design makes it difficult to react promptly to rapid contact changes, including slippage, misalignment, and external disturbances, which are pervasive in contact-rich manipulation. 

In contrast, humans perform such tasks with remarkable dexterity and robustness. Neuroscience studies~\cite{monzee2003effects, wolpert2001motor, kilteni2017sensorimotor, augurelle2003importance} suggest that this ability relies on tightly coupled multi-modal perception, predictive internal models, and rapid tactile feedback control. This combination of predictive modeling and reflexive correction is still largely missing from current visuo-tactile policies.
 

In this paper, we address these limitations from both the data and methodology perspectives. Concretely, we introduce OmniViTac, a large-scale visuo-tactile-action aligned dataset.
Building on this dataset, we further propose OmniVTA, a world-model-based visuo-tactile manipulation framework.

Specifically, OmniViTac is a large-scale visuo-tactile-action aligned dataset for contact-rich manipulation, comprising $21,879$ trajectories spanning $86$ tasks and $100+$ manipulated objects in diverse real-world scenarios. To capture the broad spectrum of contact-rich behaviors encountered in manipulation, we organize the dataset into six physics-grounded visuo-tactile interaction patterns: wiping, peeling, cutting, grasping, assembly, and in-hand adjustment.
To ensure data fidelity, OmniViTac is collected using a unified cross-embodiment platform and processed through a rigorous pipeline that preserves native sensor frequencies, performs precise temporal synchronization across vision, touch, and action streams, and incorporates automated trimming together with human-in-the-loop verification. Beyond scale and data quality, our analysis further shows that OmniViTac captures structured, pattern-specific contact dynamics and reveals two key properties of tactile signals in manipulation, namely spatial locality and contact-driven dynamics. 
These properties make OmniViTac a strong benchmark for learning visuo-tactile representations, predictive contact models, and robust manipulation policies.

Inspired by human sensorimotor control in contact-rich manipulation, an effective robotic system should possess two key capabilities~\cite{wolpert2001motor}: (1) forming feedforward anticipations of contact evolution through predictive modeling, and (2) exploiting tactile feedback for high-frequency closed-loop correction.

To this end, we propose OmniVTA, a world-model-based visuo-tactile manipulation framework composed of four tightly coupled modules. First, TactileVAE learns compact, low-dimensional tactile representations through spatio-temporal encoding and reconstructs continuous tactile deformation fields with an implicit neural decoder, providing structured tactile features for downstream prediction and policy learning. Second, the Visuo-Tactile World Model adopts a two-stream conditional generative architecture to model the temporal evolution of visual observations and tactile signals, thereby forming short-horizon feedforward predictions of contact dynamics. Third, the Adaptive Visuo-Tactile Fusion Policy leverages a Latent Tactile Differential (LTD) Encoder to capture the relationship between current tactile observations and predicted future tactile features, and adaptively fuses visual and tactile information through a contact-aware gating mechanism for action generation. Finally, the Reflexive Latent Tactile Controller uses predicted and observed tactile features to produce fine-grained corrective actions at 60,Hz, enabling rapid closed-loop correction during manipulation.

Extensive real-robot experiments across six tactile interaction pattern categories validate the effectiveness of OmniVTA from three perspectives. First, OmniVTA achieves state-of-the-art performance on diverse contact-rich manipulation tasks, with strong advantages in object diversity, geometric generalization, and disturbance robustness. The full framework further outperforms its open-loop variant, highlighting the importance of high-frequency tactile feedback for stable contact control. Second, mechanistic analyses show that the policy adaptively reweights vision and touch according to the predicted contact state, and that accurate tactile prediction is crucial for reliable action generation. Third, OmniVTA maintains stable performance under unseen geometric configurations and unseen tools, suggesting that it learns transferable contact-relevant structure and dynamics rather than memorizing task-specific trajectories. Together, these results demonstrate the benefit of predictive visuo-tactile world modeling combined with rapid tactile-feedback-based correction.

The main contributions of this paper are three-fold:
(1) We present OmniViTac, a large-scale visuo-tactile-action aligned dataset for contact-rich manipulation spanning six tactile interaction pattern categories, and reveal two key properties of tactile signals: spatial locality and contact-driven dynamics.
(2) We propose OmniVTA, a world-model-based visuo-tactile manipulation framework that integrates tactile representation learning, predictive visuo-tactile world modeling, contact-aware fusion, and high-frequency closed-loop tactile control.
(3) Extensive real-world experiments show that OmniVTA outperforms state-of-the-art baselines across diverse contact-rich manipulation tasks, with strong robustness to perturbations and effective generalization to novel configurations and objects.